% GENERATED FILE. Edit canonical lessons, then run npm run generate:books.
% canonical-source-hash: fnv1a64:c7d468fe5b2982fb
% canonical-lessons: SA-C420-prapa, SA-C420-payayati, SA-C420-mrgatrsna, SA-C420-asadayati, SA-C420-tarpayati

\chapter{Water for the Traveller}
\label{ch:sa-water-for-the-traveller}
\hlchaptermodality{420}

\begin{chapteropening}
\textbf{By the end of this chapter:} \emph{I can say where a thirsty traveller is given water.}

Complete the last lesson of chapter 420: I can say where a thirsty traveller is given water.
\end{chapteropening}

\section[\texorpdfstring{\sk{प्रपा} (prapā)}{प्रपा (prapā)}]{\sk{प्रपा} (prapā) — a roadside water shed}

\label{lesson:SA-C420-prapa}



\begin{quote}
Before the new one: say the Sanskrit for a water channel, then the Sanskrit for a puddle.
\end{quote}

\subsection*{You'll want to know: \sk{प्रपा}}
\textbf{\sk{प्रपा}} — \emph{prapā} — \textquotedblleft{}a roadside water shed\textquotedblright{}.

\begin{quote}
\sk{ग्रीष्मे} \sk{मार्गे} \sk{प्रपा} \sk{अस्ति।}
\end{quote}

\begin{quote}
\emph{grīṣme mārge prapā asti.}
\end{quote}

\begin{quote}
In summer there is a water shed by the road.
\end{quote}

A feminine noun.

\begin{cousinweb}
\textbf{\sk{प्रपा}} (\emph{prapā}) is built from \emph{pra} and the root √pā, \textquotedblleft{}to drink\textquotedblright{}, which also gives \emph{pibati}, \textquotedblleft{}drinks\textquotedblright{}; MW gives \emph{prapā} as a place for watering cattle and a shed by the road where water is given to travellers (MW).
\end{cousinweb}

\begin{grammarlens}[title={What you've built}]
One more word for getting around.
\end{grammarlens}

\subsection*{Your turn}
\begin{itemize}
\raggedright
  \item \emph{Say it:} \emph{prapā}
  \item \emph{Say it:} \emph{prapā}, once more
  \item \emph{Say it:} \emph{grīṣme mārge prapā asti.} — In summer there is a water shed by the road.
  \item \emph{Recall:} say the Sanskrit for a water channel, then the Sanskrit for a puddle, then say \emph{prapā} again
\end{itemize}

\begin{tcolorbox}[breakable,colback=teal!4,colframe=teal!35!black,title={Before you move on}]
What does \sk{प्रपा} mean? (\textquotedblleft{}A roadside water shed\textquotedblright{}.) Say it once more.
\end{tcolorbox}
\section[\texorpdfstring{\sk{पाययति} (pāyayati)}{पाययति (pāyayati)}]{\sk{पाययति} (pāyayati) — gives a drink to}

\label{lesson:SA-C420-payayati}



\begin{quote}
Before the new one: say the Sanskrit for a puddle, then the Sanskrit for a roadside water shed.
\end{quote}

\subsection*{You'll want to know: \sk{पाययति}}
\textbf{\sk{पाययति}} — \emph{pāyayati} — \textquotedblleft{}gives a drink to\textquotedblright{}.

\begin{quote}
\sk{नारी} \sk{प्रपायां} \sk{पथिकान्} \sk{जलं} \sk{पाययति।}
\end{quote}

\begin{quote}
\emph{nārī prapāyāṃ pathikān jalaṃ pāyayati.}
\end{quote}

\begin{quote}
At the water shed a woman gives the wayfarers water to drink.
\end{quote}

Past (she …): \sk{सा} \sk{पथिकान्} \sk{जलं} \sk{पायितवती}, \emph{sā pathikān jalaṃ pāyitavatī}.

\begin{cousinweb}
\textbf{\sk{पाययति}} (\emph{pāyayati}) is built on the root √pā, \textquotedblleft{}to drink\textquotedblright{}, which gives \emph{pibati}, \textquotedblleft{}drinks\textquotedblright{}; its causative, \textquotedblleft{}makes drink\textquotedblright{}, is to give a drink to, and KRM gives the causative participle \emph{pāyita} (MW).
\end{cousinweb}

\begin{grammarlens}[title={What you've built}]
One more word for getting around.
\end{grammarlens}

\subsection*{Your turn}
\begin{itemize}
\raggedright
  \item \emph{Say it:} \emph{pāyayati}
  \item \emph{Say it:} \emph{pāyayati}, once more
  \item \emph{Say it:} \emph{nārī prapāyāṃ pathikān jalaṃ pāyayati.} — At the water shed a woman gives the wayfarers water to drink.
  \item \emph{Recall:} say the Sanskrit for a puddle, then the Sanskrit for a roadside water shed, then say \emph{pāyayati} again
\end{itemize}

\begin{tcolorbox}[breakable,colback=teal!4,colframe=teal!35!black,title={Before you move on}]
What does \sk{पाययति} mean? (\textquotedblleft{}Gives a drink to\textquotedblright{}.) Say it once more.
\end{tcolorbox}
\section[\texorpdfstring{\sk{मृगतृष्णा} (mṛgatṛṣṇā)}{मृगतृष्णा (mṛgatṛṣṇā)}]{\sk{मृगतृष्णा} (mṛgatṛṣṇā) — a mirage}

\label{lesson:SA-C420-mrgatrsna}



\begin{quote}
Before the new one: say the Sanskrit for a roadside water shed, then the Sanskrit for gives a drink to.
\end{quote}

\subsection*{You'll want to know: \sk{मृगतृष्णा}}
\textbf{\sk{मृगतृष्णा}} — \emph{mṛgatṛṣṇā} — \textquotedblleft{}a mirage\textquotedblright{}.

\begin{quote}
\sk{मृगतृष्णा} \sk{जलं} \sk{न} \sk{अस्ति।}
\end{quote}

\begin{quote}
\emph{mṛgatṛṣṇā jalaṃ na asti.}
\end{quote}

\begin{quote}
A mirage is not water.
\end{quote}

A feminine noun.

\begin{cousinweb}
\textbf{\sk{मृगतृष्णा}} (\emph{mṛgatṛṣṇā}) is built from \emph{mṛga}, \textquotedblleft{}a deer, a wild animal\textquotedblright{}, and \emph{tṛṣṇā}, \textquotedblleft{}thirst\textquotedblright{}: the water a thirsty animal runs after; MW gives the compound as a mirage (MW).
\end{cousinweb}

\begin{grammarlens}[title={What you've built}]
One more word for getting around.
\end{grammarlens}

\subsection*{Your turn}
\begin{itemize}
\raggedright
  \item \emph{Say it:} \emph{mṛgatṛṣṇā}
  \item \emph{Say it:} \emph{mṛgatṛṣṇā}, once more
  \item \emph{Say it:} \emph{mṛgatṛṣṇā jalaṃ na asti.} — A mirage is not water.
  \item \emph{Recall:} say the Sanskrit for a roadside water shed, then the Sanskrit for gives a drink to, then say \emph{mṛgatṛṣṇā} again
\end{itemize}

\begin{tcolorbox}[breakable,colback=teal!4,colframe=teal!35!black,title={Before you move on}]
What does \sk{मृगतृष्णा} mean? (\textquotedblleft{}A mirage\textquotedblright{}.) Say it once more.
\end{tcolorbox}
\section[\texorpdfstring{\sk{आसादयति} (āsādayati)}{आसादयति (āsādayati)}]{\sk{आसादयति} (āsādayati) — comes upon; reaches}

\label{lesson:SA-C420-asadayati}



\begin{quote}
Before the new one: say the Sanskrit for gives a drink to, then the Sanskrit for a mirage.
\end{quote}

\subsection*{You'll want to know: \sk{आसादयति}}
\textbf{\sk{आसादयति}} — \emph{āsādayati} — \textquotedblleft{}comes upon; reaches\textquotedblright{}.

\begin{quote}
\sk{सायं} \sk{पथिकः} \sk{प्रपाम्} \sk{आसादयति।}
\end{quote}

\begin{quote}
\emph{sāyaṃ pathikaḥ prapām āsādayati.}
\end{quote}

\begin{quote}
In the evening the wayfarer comes upon a water shed.
\end{quote}

Past (he …): \sk{सः} \sk{प्रपाम्} \sk{आसादितवान्}, \emph{saḥ prapām āsāditavān}.

\begin{cousinweb}
\textbf{\sk{आसादयति}} (\emph{āsādayati}) is built from \emph{ā} and the causative of the root √sad, \textquotedblleft{}to sit\textquotedblright{}; MW gives its participle \emph{āsādita} as reached (MW).
\end{cousinweb}

\begin{grammarlens}[title={What you've built}]
One more word for getting around.
\end{grammarlens}

\subsection*{Your turn}
\begin{itemize}
\raggedright
  \item \emph{Say it:} \emph{āsādayati}
  \item \emph{Say it:} \emph{āsādayati}, once more
  \item \emph{Say it:} \emph{sāyaṃ pathikaḥ prapām āsādayati.} — In the evening the wayfarer comes upon a water shed.
  \item \emph{Recall:} say the Sanskrit for gives a drink to, then the Sanskrit for a mirage, then say \emph{āsādayati} again
\end{itemize}

\begin{tcolorbox}[breakable,colback=teal!4,colframe=teal!35!black,title={Before you move on}]
What does \sk{आसादयति} mean? (\textquotedblleft{}Comes upon; reaches\textquotedblright{}.) Say it once more.
\end{tcolorbox}
\section[\texorpdfstring{\sk{तर्पयति} (tarpayati)}{तर्पयति (tarpayati)}]{\sk{तर्पयति} (tarpayati) — refreshes; satisfies}

\label{lesson:SA-C420-tarpayati}



\begin{quote}
Before the new one: say the Sanskrit for a mirage, then the Sanskrit for comes upon.
\end{quote}

\subsection*{You'll want to know: \sk{तर्पयति}}
\textbf{\sk{तर्पयति}} — \emph{tarpayati} — \textquotedblleft{}refreshes; satisfies\textquotedblright{}.

\begin{quote}
\sk{शीतलं} \sk{जलं} \sk{श्रान्तं} \sk{पथिकं} \sk{तर्पयति।}
\end{quote}

\begin{quote}
\emph{śītalaṃ jalaṃ śrāntaṃ pathikaṃ tarpayati.}
\end{quote}

\begin{quote}
Cool water refreshes the tired wayfarer.
\end{quote}

Past (it …): \sk{जलं} \sk{पथिकं} \sk{तर्पितवत्}, \emph{jalaṃ pathikaṃ tarpitavat}.

\begin{cousinweb}
\textbf{\sk{तर्पयति}} (\emph{tarpayati}) is built on the root √tṛp, \textquotedblleft{}to be satisfied\textquotedblright{}, in its causative, \textquotedblleft{}to satisfy, to refresh\textquotedblright{}; MW gives the participle \emph{tarpita}, \textquotedblleft{}satisfied\textquotedblright{} (MW).
\end{cousinweb}

\begin{grammarlens}[title={What you've built}]
One more word for getting around.
\end{grammarlens}

\subsection*{Your turn}
\begin{itemize}
\raggedright
  \item \emph{Say it:} \emph{tarpayati}
  \item \emph{Say it:} \emph{tarpayati}, once more
  \item \emph{Say it:} \emph{śītalaṃ jalaṃ śrāntaṃ pathikaṃ tarpayati.} — Cool water refreshes the tired wayfarer.
  \item \emph{Recall:} say the Sanskrit for a mirage, then the Sanskrit for comes upon, then say \emph{tarpayati} again
\end{itemize}

\begin{tcolorbox}[breakable,colback=teal!4,colframe=teal!35!black,title={Before you move on}]
What does \sk{तर्पयति} mean? (\textquotedblleft{}Refreshes; satisfies\textquotedblright{}.) Say it once more.
\end{tcolorbox}
